\section{Responsibility and duty: you have the power to prevent this}
\label{sec:responsibility}

The fifth appeal is addressed to the legislator directly. It is not sympathy but
obligation. You have the power to prevent this, and future generations will judge
what we do now. This is the moment the narrative turns from describing a problem to
naming the body that can solve it and the precise instrument, H. R. 9510, that
lets it. Duty creates a sense of agency that sympathy alone does not.

Responsibility is exercised through a known sequence of actions. \emph{Advocacy}
speaks in favor of safe medical AI. \emph{Coalition building} assembles patient
groups, hospital networks, technology developers, and medical boards so the bill
visibly carries wide support. \emph{Policy entrepreneurship} matches the
documented problem of unverified AI to this specific legislative fix. And when two
chambers pass different texts, \emph{reconciliation} resolves them into one.
Credible \emph{testimony} at \emph{committee hearings} does not always change a
vote, but it reliably informs the record, shapes the floor debate, and clarifies
the language a committee writes during \emph{markup} \cite{moreland2015hearing}.
\autoref{fig:duty} shows the path on which these actions operate.

\begin{narrativefig}
\node[nfone] (a) at (0,0) {Introduce};
\node[nftwo] (b) at (3.2,0) {Hearings and testimony};
\node[nfact] (c) at (6.7,0) {Markup};
\node[nftwo] (d) at (10.0,0) {Floor passage};
\node[nfthree] (e) at (5.0,-2.1) {Reconciliation};
\node[nfhope] (f) at (8.5,-2.1) {Enacted law};
\draw[nfedge] (a)--(b); \draw[nfedge] (b)--(c); \draw[nfedge] (c)--(d);
\draw[nfedge] (d) to[out=-90,in=0] (e.east);
\draw[nfedge] (e)--(f);
\draw[nfedged] (c) to[out=-120,in=120] (b.south);
\end{narrativefig}
\vspace{-0.15cm}
\figcaption{Figure 7. The duty to move a bill from introduction to enactment
(Mermaid 03 and 14 as colored TikZ), with testimony feeding markup and
reconciliation joining two chambers.}

\autoref{tab:duty} maps each main action to the committee step where it has the
most effect, so a member can see exactly where to apply effort.

\begin{table}[H]
\footnotesize
\begin{tabularx}{\textwidth}{@{}L{4.0cm} Y@{}}
\toprule
\textbf{Main action} & \textbf{Where it has effect}\\
\midrule
Advocacy & Building awareness before referral\\
Coalition building & Demonstrating support at hearings\\
Policy entrepreneurship & Framing the problem during markup\\
Reconciliation & Resolving House and Senate texts\\
\bottomrule
\end{tabularx}
\caption{Table 5. Each main action and the legislative step it serves.}
\label{tab:duty}
\end{table}

Responsibility is the bridge from feeling to action. It tells the committee that
the next move is theirs and that the record will remember it. The quiet question
becomes a charge. If the power to prevent foreseeable harm rests here, why would we
leave it unused?
